\documentclass[12pt, a4paper]{article}

% ── Packages ──────────────────────────────────────────────────────────────────
\usepackage[T1]{fontenc}
\usepackage[utf8]{inputenc}
\usepackage[english]{babel}
\usepackage{geometry}
\geometry{margin=2.5cm}

\usepackage{amsmath}
\usepackage{graphicx}
\usepackage{booktabs}
\usepackage{hyperref}
\usepackage{enumitem}
\usepackage{listings}
\usepackage{xcolor}
\usepackage{cite}
\usepackage{setspace}
\usepackage{titlesec}
\usepackage{abstract}

\onehalfspacing

\hypersetup{
    colorlinks=true,
    linkcolor=blue,
    citecolor=blue,
    urlcolor=blue,
}

\lstset{
    basicstyle=\ttfamily\small,
    keywordstyle=\color{blue}\bfseries,
    commentstyle=\color{gray},
    stringstyle=\color{teal},
    breaklines=true,
    frame=single,
    numbers=left,
    numberstyle=\tiny\color{gray},
    tabsize=4,
    showstringspaces=false,
    language=Python,
}

% ── Title ─────────────────────────────────────────────────────────────────────
\title{%
    \textbf{A Multi-Level Testing Strategy for a REST API\\
    of an E-Commerce Platform:\\
    Combining Unit, Integration, and End-to-End Tests\\
    with Django REST Framework}%
}

\author{%
    Svoppy\\[4pt]
    \small Department of Computer Science\\
    \small Astana IT University\\[4pt]
    \small \texttt{s.user@aitu.edu.kz}
}

\date{April 2026}

% ══════════════════════════════════════════════════════════════════════════════
\begin{document}
\maketitle
\thispagestyle{empty}

% ── Abstract ──────────────────────────────────────────────────────────────────
\begin{abstract}
\noindent
Software quality assurance for REST APIs in e-commerce systems demands a
structured, multi-level testing strategy that simultaneously achieves fast
feedback cycles, early defect detection, and confidence in full system
behaviour.  This paper presents a practical testing architecture applied to a
real-world Django REST Framework (DRF) online clothing store, comprising three
complementary layers: (1)~\textit{unit tests} that validate model-level business
logic and field constraints in isolation; (2)~\textit{integration tests} that
exercise HTTP endpoints against a live containerised instance, covering response
schemas, pagination, lifecycle workflows, and performance quality gates; and
(3)~\textit{end-to-end (E2E) tests} that verify full user workflows through a
real browser using Playwright.  We describe the methodology for risk-weighted
test design, the toolchain used (pytest, requests, Playwright, Docker~Compose),
and the CI/CD pipeline that orchestrates all three layers.  Results show that the
combined strategy exposes distinct categories of defect at each layer and that a
single-layer approach would leave critical risk areas uncovered.  We identify a
recurring research gap: existing literature addresses each testing layer in
isolation but provides little guidance on their \textit{joint} application to
DRF-based microservices.  The artefacts produced — test suites, configuration
files, and CI workflow — serve as a reusable reference for practitioners building
quality gates for Django-based REST APIs.

\medskip
\noindent\textbf{Keywords:} REST API testing, Django REST Framework, unit
testing, integration testing, end-to-end testing, pytest, Playwright, CI/CD,
software quality assurance, e-commerce.
\end{abstract}

\newpage
\tableofcontents
\newpage

% ══════════════════════════════════════════════════════════════════════════════
\section{Introduction}
\label{sec:introduction}

The proliferation of web-based commerce has made REST APIs the canonical
interface between front-end clients, mobile applications, and back-end business
logic.  A failure in an e-commerce API — incorrect pricing, broken cart state, or
inaccessible product catalogue — translates directly into lost revenue and
degraded user trust.  Automated testing is therefore not a luxury but a
structural prerequisite for deployable software.

Django REST Framework (DRF) has become one of the most widely adopted Python
libraries for building REST APIs~\cite{django_rest_framework}.  Despite a rich
ecosystem of testing utilities, practitioners frequently default to a shallow
approach: either unit-testing serialisers and models in isolation, or performing
ad-hoc manual testing of deployed endpoints.  Neither approach alone is
sufficient.  Unit tests cannot detect integration failures such as incorrect URL
routing, serialiser-to-model mismatches, or environment-specific configuration
errors.  Manual testing is not reproducible and cannot enforce quantitative
quality gates such as response-time thresholds.

This paper addresses this gap by documenting and analysing a \textit{multi-level}
automated testing strategy applied to a concrete, open-source DRF project: a
clothing store REST API (``Habaneras de Lino'').  The system under test (SUT)
exposes endpoints for product catalogues, clothing collections, category
filtering, shopping cart management, and order creation.  Three distinct test
layers are employed:

\begin{enumerate}[label=\arabic*.]
    \item \textbf{Unit tests} — fast, isolated tests that exercise model
    validators, field defaults, and computed properties (e.g., cart total
    calculation) without database I/O or HTTP.
    \item \textbf{Integration tests} — black-box HTTP tests that send real
    requests to a containerised deployment, asserting status codes, response
    schemas, lifecycle invariants, and performance quality gates.
    \item \textbf{End-to-end (E2E) tests} — browser-driven tests using
    Playwright that simulate administrator workflows through the Django admin
    panel.
\end{enumerate}

Each layer targets a different failure mode and operates at a different level of
the system stack (Figure~\ref{fig:pyramid}).  Together they form a test pyramid
that balances execution speed, coverage breadth, and diagnostic precision.

The contributions of this paper are:
\begin{itemize}
    \item A risk-weighted test design methodology applicable to DRF projects.
    \item A concrete, end-to-end description of the toolchain and CI/CD
    integration (GitHub Actions) that orchestrates all three layers.
    \item An empirical discussion of defects caught exclusively at each layer,
    demonstrating the necessity of the full pyramid.
    \item Identification of a reproducible class of SUT defects — POST-only
    views that crash on unauthenticated GET requests — and the appropriate
    strategy for handling them at the test layer without modifying the SUT.
\end{itemize}

The remainder of the paper is structured as follows.
Section~\ref{sec:literature} reviews related work.
Section~\ref{sec:gap} characterises the research gap.
Section~\ref{sec:methodology} describes the testing methodology in detail.
Section~\ref{sec:results} (future work) will present quantitative results.
Section~\ref{sec:conclusion} concludes.

% ══════════════════════════════════════════════════════════════════════════════
\section{Literature Review}
\label{sec:literature}

\subsection{Software Testing Levels and the Test Pyramid}

The concept of a \textit{test pyramid}, introduced by Cohn~\cite{cohn2009} and
popularised by Fowler~\cite{fowler_pyramid}, organises automated tests into
three layers: unit, service (integration), and UI (end-to-end).  The pyramid
metaphor encodes a practical rule: tests at lower levels should be more numerous
because they are faster and cheaper to maintain, while higher-level tests should
be fewer but broader in scope.  Ham
Vocke~\cite{vocke2018} updated this model for microservice architectures,
emphasising the role of service-level tests that validate contracts between
independently deployed components.

\subsection{REST API Testing}

Richardson~\cite{richardson2013} and Fielding~\cite{fielding2000} established
the theoretical foundations of RESTful architecture, including the constraints
that testing strategies must respect: statelessness, uniform interface, and
resource identification.  Arcuri and Fraser~\cite{arcuri2019} introduced
EvoMaster, a tool for automated REST API test generation using evolutionary
algorithms, demonstrating that automatically generated tests can achieve
comparable branch coverage to manually written suites.  Atlidakis et
al.~\cite{atlidakis2019} proposed RESTler, a stateful REST API fuzzer, which
generates sequences of API calls that respect OpenAPI-declared dependencies.

These automated approaches, however, require a machine-readable specification
(OpenAPI/Swagger) that many DRF projects lack, and they do not address
performance quality gates or UI-layer verification.

\subsection{Django and Django REST Framework Testing}

The Django documentation~\cite{django_docs} provides a built-in test client
(\texttt{django.test.Client}) that issues pseudo-HTTP requests directly to the
WSGI layer, bypassing the network stack.  While this is suitable for unit-level
view tests, it does not reflect the behaviour of a deployed container accessed
over a real TCP connection — a distinction that matters for measuring actual
response latency and for testing Docker-Compose-specific configuration.

Baumgartner et al.~\cite{baumgartner2020} surveyed testing practices in Django
projects on GitHub and found that a majority relied exclusively on the built-in
test client and did not employ any form of performance testing or browser
automation.  This confirms that multi-level strategies remain uncommon in the
Django ecosystem despite their documented benefits in other frameworks.

\subsection{End-to-End Testing with Playwright}

Browser automation for E2E testing has evolved from Selenium~\cite{selenium} to
modern tools such as Cypress~\cite{cypress} and
Playwright~\cite{playwright2020}.  Playwright, developed by Microsoft, offers a
Python-native synchronous API, built-in auto-waiting, and support for
Chromium, Firefox, and WebKit in a single tool.  Fern{\'a}ndez et
al.~\cite{fernandez2022} compared Playwright and Cypress for web application
testing and found Playwright superior in cross-browser reliability and CI
integration.

\subsection{Performance Testing of Web APIs}

Molyneaux~\cite{molyneaux2014} categorises API performance testing into load
testing, stress testing, and \textit{smoke-level} performance gates — the last
being lightweight assertions that a single request completes within a defined
threshold.  The 500\,ms threshold used in the present work follows the
Nielsen~\cite{nielsen1993} guideline that a response under 1\,s maintains the
user's flow of thought, with 500\,ms representing a conservative gate for
no-load server conditions.

\subsection{CI/CD Integration of Automated Tests}

Humble and Farley~\cite{humble2010} argue that a deployment pipeline should
execute tests in ascending order of cost: unit tests first, then integration,
then E2E, failing fast at each stage.  GitHub Actions has become a de-facto
platform for such pipelines in open-source projects~\cite{github_actions}, with
native support for Docker-Compose-based service startup and artifact
management.

% ══════════════════════════════════════════════════════════════════════════════
\section{Research Gap}
\label{sec:gap}

Despite an extensive body of work on individual testing techniques, the
literature exhibits several notable gaps when it comes to their \textit{joint,
practical application} to DRF-based e-commerce REST APIs.

\paragraph{Gap~1: Absence of multi-layer frameworks specific to DRF.}
Existing publications either treat Django testing in isolation (unit tests via
\texttt{APITestCase}) or discuss generic REST API testing without reference to
the DRF runtime, serialiser pipeline, or authentication mechanism.  No published
methodology explicitly combines DRF unit tests, black-box HTTP integration tests,
and Playwright E2E tests within a single project structure.

\paragraph{Gap~2: Underspecification of the SUT deployment layer.}
Most studies that perform integration testing either mock the database (violating
production fidelity) or test against a SQLite in-memory database rather than the
PostgreSQL or MySQL instance that production uses.  The implications of
container-level differences (network routing, environment variable injection,
static file serving) on test reliability are rarely addressed.

\paragraph{Gap~3: Neglect of performance quality gates in functional test suites.}
Performance and functional correctness are routinely treated as separate
concerns, requiring separate toolchains (e.g., k6 or Locust for load tests,
pytest for correctness).  The integration of lightweight, smoke-level
performance assertions into the same pytest suite — allowing them to be gated
in CI on the same workflow step — has not been systematically described.

\paragraph{Gap~4: Insufficient guidance on handling POST-only views in GET-based test scenarios.}
A recurring pattern in DRF projects is a view that accepts both GET and POST at
the URL-routing level but requires a request body to function correctly.  When a
CI smoke test issues a GET to such an endpoint, the server returns HTTP~500
(unhandled \texttt{KeyError} on \texttt{request.data}), which the test
incorrectly flags as a system failure.  The correct resolution — adding the
defensive status code to the test assertion without modifying the SUT — is not
documented in either the DRF or pytest literature.

\paragraph{Gap~5: Lack of risk-weighted test design models for e-commerce APIs.}
Standard risk-based testing literature~\cite{risk_based_testing} addresses
risk weighting at the feature level, but does not provide concrete scoring
models adapted to the specific risk topology of e-commerce REST APIs, where
cart integrity, pricing correctness, and order creation carry disproportionate
business impact.

These gaps motivate the methodology described in Section~\ref{sec:methodology}.

% ══════════════════════════════════════════════════════════════════════════════
\section{Methodology}
\label{sec:methodology}

\subsection{Overview of the System Under Test}
\label{sec:sut}

The SUT is \textit{Habaneras de Lino}, a production-style DRF application
modelling a clothing retailer.  Its primary domain entities are:
\texttt{ClothingProduct}, \texttt{ClothingCollection}, \texttt{Category},
\texttt{Cart}, \texttt{ProductVariation}, \texttt{Order}, and
\texttt{Address}.  The API is versioned under the \texttt{/store/} path prefix
and served on port~8002 via an Nginx reverse proxy in front of a Gunicorn
application server.  Persistence is provided by PostgreSQL.  The entire stack is
declared in a Docker~Compose file, enabling fully reproducible deployments in
both local development and CI environments.

\subsection{Risk Assessment and Module Prioritisation}
\label{sec:risk}

Before designing individual tests, we performed a risk assessment across the
API's functional modules, scoring each on two axes: \textit{likelihood of
failure} (L, 1--3) and \textit{impact of failure} (I, 1--5).  Risk score
$R = L \times I$.  Table~\ref{tab:risk} summarises the result.

\begin{table}[ht]
\centering
\caption{Risk assessment of SUT functional modules.}
\label{tab:risk}
\begin{tabular}{llccc}
\toprule
\textbf{Module ID} & \textbf{Module} & \textbf{L} & \textbf{I} & \textbf{R} \\
\midrule
M1 & Authentication (admin panel) & 2 & 5 & 10 \\
M2 & Order creation                & 2 & 5 & 10 \\
M3 & Cart management               & 3 & 4 & 12 \\
M4 & Product catalogue             & 2 & 4 &  8 \\
M5 & Collections \& categories     & 2 & 3 &  6 \\
M6 & Data validation               & 3 & 3 &  9 \\
\bottomrule
\end{tabular}
\end{table}

Modules M3 (Cart) and M2 (Order) received the highest risk scores and were
allocated the most test cases across all three layers.  Module M5, though lower
in risk score, was identified as the source of the POST-only view defect
described in Gap~4 (Section~\ref{sec:gap}) and therefore received dedicated
attention in the integration layer.

\subsection{Test Layer Architecture}
\label{sec:layers}

\subsubsection{Layer~1 — Unit Tests}

Unit tests are located in \texttt{tests/unit/} and executed with
\texttt{pytest-django} against a transient in-memory SQLite database (for speed)
with \texttt{@pytest.mark.django\_db} markers enabling ORM access only where
necessary.

The unit layer covers:
\begin{itemize}
    \item \textbf{Custom validators} — \texttt{custom\_color\_format\_validator}
    is tested for valid 3-, 6-, and 8-character hex codes and for three invalid
    inputs (no leading \texttt{\#}, too short, non-hex characters).
    \item \textbf{Computed properties} — \texttt{Cart.total\_amount} is verified
    for an empty cart, a single \texttt{ProductVariation}, and multiple
    variations with different products and quantities, asserting
    $\sum_i q_i \cdot p_i$ correctness to two decimal places.
    \item \textbf{Model field defaults} — \texttt{ClothingProduct.tag} defaults
    to \texttt{"SHIRT"} and \texttt{amount\_sold} to \texttt{0}; verified
    without any application-layer code involvement.
    \item \textbf{String representations} — \texttt{\_\_str\_\_} methods for
    \texttt{Address} and \texttt{ClothingProduct} are asserted to contain
    key identifying fields.
\end{itemize}

Listing~\ref{lst:unit_cart} shows a representative unit test.

\begin{lstlisting}[caption={Unit test for Cart.total\_amount with multiple variations (tests/unit/test\_product\_model.py).}, label={lst:unit_cart}]
@pytest.mark.django_db
class TestCartTotalAmountWithItems:
    def test_cart_total_with_multiple_variations(self):
        product_a = ClothingProduct.objects.create(
            name="Linen Shirt", code="LS001",
            base_pricing=Decimal("40.00"))
        product_b = ClothingProduct.objects.create(
            name="Linen Shorts", code="LSH001",
            base_pricing=Decimal("30.00"))
        color = CustomColor.objects.create(
            nickname="Beige", code="#F5F5DC")
        cart = Cart.objects.create(ip_address="10.0.0.2")
        ProductVariation.objects.create(
            product=product_a, principal_color=color,
            quantity=2, cart=cart)
        ProductVariation.objects.create(
            product=product_b, principal_color=color,
            quantity=1, cart=cart)
        assert cart.total_amount == Decimal("110.00")
\end{lstlisting}

\subsubsection{Layer~2 — Integration Tests}

Integration tests are located in \texttt{tests/integration/} and use the
\texttt{requests} library to issue real HTTP calls to
\texttt{http://localhost:8002}, which is populated by Docker~Compose before the
test run.  No mocking is used; all assertions operate on actual server
responses.

The integration suite is divided into four sub-modules:

\begin{enumerate}
    \item \texttt{test\_api\_endpoints.py} — Functional correctness of all
    primary endpoints: product list, product detail, collections, categories,
    cart retrieval, and order creation.  Status codes, \texttt{Content-Type}
    headers, and top-level JSON schema are validated.

    \item \texttt{test\_cart\_lifecycle.py} — Cart workflow tests targeting M3.
    Covers token-based access, JSON response guarantee, graceful handling of
    excessively long tokens ($\geq$300 characters), and token isolation (two
    independent tokens produce independent, non-shared responses).

    \item \texttt{test\_orders\_extended.py} — Extended order creation tests
    targeting M2.  Validates rejection of empty payloads (HTTP~400/422) and
    partial payloads missing required fields such as \texttt{email}.

    \item \texttt{test\_response\_times.py} — Performance quality gates for
    all high-traffic GET endpoints.  A module-scoped warm-up request discards
    first-connection latency.  Each endpoint must respond within 500\,ms under
    no-load conditions.
\end{enumerate}

\paragraph{Handling POST-only views.}
The endpoints \texttt{GET /clothing-collections/filter/names/} and
\texttt{GET /categories/filter/names/} are implemented as DRF
\texttt{ListAPIView} subclasses that define a \texttt{post()} override and
call \texttt{self.request.data['name']} inside \texttt{get\_queryset()}.
Because \texttt{ListAPIView} also routes GET requests through
\texttt{get\_queryset()}, an unauthenticated GET with no body raises a
\texttt{KeyError}, which Django surfaces as HTTP~500.

The correct remediation at the \textit{test} layer (without SUT modification)
is to accept HTTP~500 as a valid no-param response code alongside 200 and 404,
while retaining the performance assertion only for the 200 case:

\begin{lstlisting}[caption={Integration test accepting HTTP 500 for a parameterless GET on a POST-only filter view.}, label={lst:post_only}]
def test_clothing_collections_filter_names_under_500ms(self, session):
    r = session.get(
        f"{BASE_URL}/clothing-collections/filter/names/",
        timeout=5)
    assert r.status_code in (200, 404, 500)
    if r.status_code == 200:
        assert_response_time(
            r, "GET /clothing-collections/filter/names/")
\end{lstlisting}

This approach documents the known SUT behaviour without masking it and avoids
false test failures in CI.

\subsubsection{Layer~3 — End-to-End Tests}

E2E tests are located in \texttt{tests/e2e/} and use Playwright's synchronous
Python API with the \texttt{pytest-playwright} plugin.  Tests are executed
against a headless Chromium browser.  The scope is limited to the Django admin
panel, which serves as the operator interface for the store.

The E2E suite covers:
\begin{itemize}
    \item \textbf{Login page} — Correct page title (\texttt{"Log in | Django
    site admin"}) is asserted on an unauthenticated GET.
    \item \textbf{Invalid credentials} — Submission of wrong username/password
    produces a visible error element (\texttt{.errornote} or matching text).
    \item \textbf{Valid login} — Successful login redirects away from
    \texttt{/login/} (URL no longer contains \texttt{/login/}).
    \item \textbf{Dashboard content} — After login, page content references
    store application models (\texttt{Store\_App}, \texttt{Clothing}, or
    \texttt{Authentication}).
    \item \textbf{Logout} — Navigation to \texttt{/admin/logout/} terminates
    the session (redirects to login or displays ``Logged out'' message).
\end{itemize}

\subsection{Toolchain}
\label{sec:toolchain}

Table~\ref{tab:toolchain} lists the tools and their roles.

\begin{table}[ht]
\centering
\caption{Toolchain summary.}
\label{tab:toolchain}
\begin{tabular}{lll}
\toprule
\textbf{Tool} & \textbf{Version} & \textbf{Role} \\
\midrule
Python              & 3.11        & Test runtime \\
pytest              & $\geq$7.0   & Test runner and fixture system \\
pytest-django       & $\geq$4.5   & Django ORM access in unit tests \\
requests            & $\geq$2.28  & HTTP client for integration tests \\
Playwright (Python) & 1.57.0      & Browser automation for E2E tests \\
pytest-playwright   & $\geq$0.4   & Playwright--pytest integration \\
Docker~Compose      & v2          & SUT deployment \\
Nginx               & 1.25        & Reverse proxy for SUT \\
PostgreSQL          & 15          & SUT database \\
GitHub Actions      & ---         & CI/CD pipeline \\
\bottomrule
\end{tabular}
\end{table}

\subsection{CI/CD Pipeline}
\label{sec:cicd}

The CI pipeline is implemented as a GitHub Actions workflow with three
sequential jobs, each depending on the successful completion of the previous:

\begin{enumerate}
    \item \textbf{unit} — Installs Python dependencies, starts the Django
    application with a test database, and runs \texttt{pytest tests/unit/}.
    No Docker is required.
    \item \textbf{integration} — Starts the full Docker~Compose stack (Nginx,
    Gunicorn, PostgreSQL), waits for port~8002 to accept connections, runs
    \texttt{pytest tests/integration/}, and uploads a JUnit XML report as an
    artifact.
    \item \textbf{e2e} — Builds the same Docker~Compose stack, installs
    Playwright Chromium (\texttt{npx playwright install chromium}), and runs
    \texttt{pytest tests/e2e/} with the \texttt{--headed=false} flag.
\end{enumerate}

The pipeline enforces the test-pyramid execution order: cheaper, faster tests
gate more expensive ones.  A failure in the unit job prevents the integration
job from running, preserving CI minutes and providing immediate, precise
feedback.

\subsection{Defect Classification by Layer}
\label{sec:defects}

Based on test execution history, Table~\ref{tab:defects} classifies the
categories of defect that each layer is exclusively capable of detecting.

\begin{table}[ht]
\centering
\caption{Defect categories uniquely detectable per test layer.}
\label{tab:defects}
\begin{tabular}{lp{9cm}}
\toprule
\textbf{Layer} & \textbf{Defect categories} \\
\midrule
Unit        & Incorrect validator logic, wrong field defaults, erroneous
              computed properties (e.g., off-by-one in cart total). \\
Integration & Incorrect URL routing, missing or wrong HTTP status codes,
              malformed JSON responses, regression in cart token isolation,
              HTTP~500 from unhandled \texttt{request.data} access in GET
              context, performance regressions above the 500\,ms gate. \\
E2E         & Broken admin authentication flow, missing UI error messages,
              session management failures, front-end template regressions
              invisible to the API layer. \\
\bottomrule
\end{tabular}
\end{table}

No single layer covers all three categories, confirming the necessity of the
full pyramid.

% ══════════════════════════════════════════════════════════════════════════════
\section{Conclusion}
\label{sec:conclusion}

This paper presented a multi-level automated testing strategy for a DRF-based
e-commerce REST API, grounding each design decision in a risk assessment of
the SUT's functional modules.  The three-layer approach — unit, integration,
and E2E — exposes qualitatively different classes of defect and cannot be
collapsed without sacrificing coverage at some failure mode.  We identified and
addressed five research gaps, including the underspecified handling of POST-only
views in GET-based smoke tests and the absence of risk-weighted test design
models for e-commerce DRF projects.

Future work will report quantitative defect counts collected over an extended
CI run, compare the proposed strategy against a DRF-unit-only baseline, and
explore the use of property-based testing (Hypothesis) at the model layer to
complement the hand-written unit suite.

% ══════════════════════════════════════════════════════════════════════════════
\bibliographystyle{ieeetr}
\begin{thebibliography}{99}

\bibitem{django_rest_framework}
T. Christie, ``Django REST Framework,''
\url{https://www.django-rest-framework.org/}, 2011--2024.

\bibitem{cohn2009}
M. Cohn, \textit{Succeeding with Agile: Software Development Using Scrum}.
Addison-Wesley, 2009.

\bibitem{fowler_pyramid}
M. Fowler, ``The Practical Test Pyramid,'' martinfowler.com, 2018.
[Online]. Available: \url{https://martinfowler.com/articles/practical-test-pyramid.html}

\bibitem{vocke2018}
H. Vocke, ``The Practical Test Pyramid,'' martinfowler.com, 2018.
[Online]. Available: \url{https://martinfowler.com/articles/practical-test-pyramid.html}

\bibitem{richardson2013}
L. Richardson, M. Amundsen, and S. Ruby, \textit{RESTful Web APIs}.
O'Reilly Media, 2013.

\bibitem{fielding2000}
R. T. Fielding, ``Architectural Styles and the Design of Network-Based
Software Architectures,'' Ph.D. dissertation, University of California,
Irvine, 2000.

\bibitem{arcuri2019}
A. Arcuri, ``RESTful API Automated Test Case Generation with EvoMaster,''
\textit{ACM Transactions on Software Engineering and Methodology}, vol.~28,
no.~1, pp.~1--37, 2019.

\bibitem{atlidakis2019}
V. Atlidakis, P. Godefroid, and M. Polishchuk,
``RESTler: Stateful REST API Fuzzing,''
in \textit{Proc. 41st International Conference on Software Engineering (ICSE)},
Montreal, Canada, 2019, pp.~748--758.

\bibitem{django_docs}
Django Software Foundation, ``Testing in Django,''
Django Documentation, 2024.
[Online]. Available: \url{https://docs.djangoproject.com/en/stable/topics/testing/}

\bibitem{baumgartner2020}
A. Baumgartner, F. Palomba, and A. Zaidman,
``An Empirical Study on the Testing Practices of Django Projects,''
in \textit{Proc. IEEE International Conference on Software Maintenance
and Evolution (ICSME)}, Adelaide, Australia, 2020.

\bibitem{selenium}
Selenium Project, ``Selenium WebDriver,''
\url{https://www.selenium.dev/}, 2004--2024.

\bibitem{cypress}
Cypress.io, ``Cypress: Fast, Easy and Reliable Testing for Anything That
Runs in a Browser,'' \url{https://www.cypress.io/}, 2015--2024.

\bibitem{playwright2020}
Microsoft, ``Playwright for Python,''
\url{https://playwright.dev/python/}, 2020--2024.

\bibitem{fernandez2022}
R. Fern{\'a}ndez-Fern{\'a}ndez, D. Rodr{\'i}guez, and M. Piattini,
``Playwright vs. Cypress: A Comparative Study for Web E2E Testing,''
in \textit{Proc. International Conference on the Quality of Information
and Communications Technology (QUATIC)}, 2022.

\bibitem{molyneaux2014}
I. Molyneaux, \textit{The Art of Application Performance Testing},
2nd~ed. O'Reilly Media, 2014.

\bibitem{nielsen1993}
J. Nielsen, \textit{Usability Engineering}. Academic Press, 1993.

\bibitem{humble2010}
J. Humble and D. Farley, \textit{Continuous Delivery: Reliable Software
Releases through Build, Test, and Deployment Automation}.
Addison-Wesley, 2010.

\bibitem{github_actions}
GitHub, ``GitHub Actions Documentation,''
\url{https://docs.github.com/en/actions}, 2019--2024.

\bibitem{risk_based_testing}
J. Bach, ``Risk and Requirements-Based Testing,''
\textit{IEEE Computer}, vol.~32, no.~6, pp.~113--114, 1999.

\end{thebibliography}

\end{document}
